% !TEX root = ../thesis.tex

\chapter{Používateľská príručka}\label{app:user.manual}

Táto príručka opisuje používanie aplikácie Energy AI Assistant. Zameriava sa na inštaláciu, spustenie a ovládanie rozhrania.

\section{Požiadavky}

Na spustenie aplikácie sú potrebné tieto komponenty:
\begin{itemize}
	\item Linux alebo macOS. Windows je možný cez WSL.
	\item Python 3.10 alebo novší a nástroj \texttt{pip}.
	\item Voľne dostupný priestor pre modely (niekoľko GB).
	\item Internet na prvé stiahnutie základných modelov z Hugging Face.
	\item GPU NVIDIA s CUDA pre rozumnú rýchlosť. Bez GPU je odozva pomalá.
\end{itemize}

Pri niektorých modeloch (napr. Gemma) môže byť potrebné schválenie licencie v~Hugging Face. V~takom prípade je nutné použiť osobný prístupový token.

\section{Inštalácia}

Postup je jednoduchý. Príklady sú uvedené pre Linux.

\begin{enumerate}
	\item Vytvorte virtuálne prostredie:
\begin{verbatim}
python3 -m venv venv
source venv/bin/activate
\end{verbatim}

	\item Nainštalujte balíky:
\begin{verbatim}
pip install -r requirements.txt
\end{verbatim}

	\item Voliteľne vytvorte súbor \texttt{.env}:
\begin{verbatim}
cp .env.example .env
\end{verbatim}
\end{enumerate}

Ak používate GPU, odporúča sa mať nainštalovanú CUDA verziu PyTorch. Presný príkaz závisí od verzie CUDA.

\section{Spustenie aplikácie}

Aplikáciu je možné spustiť viacerými spôsobmi.

\textbf{Spustenie cez hlavný súbor:}
\begin{verbatim}
python energy_chat/main.py
\end{verbatim}

\textbf{Spustenie cez Uvicorn:}
\begin{verbatim}
python -m uvicorn energy_chat.main:app --host 0.0.0.0 --port 8000
\end{verbatim}

\textbf{Spustenie cez skript:}
\begin{verbatim}
./start_server.sh
\end{verbatim}

Po spustení sa v~konzole zobrazí adresa servera. Predvolený port je nastavený v~\path{energy_chat/core/config.py}. V~aktuálnej konfigurácii je to port 8002. V~prehliadači otvorte adresu \texttt{http://localhost:8002} alebo port, ktorý ste nastavili.

\section{Použitie aplikácie}

Rozhranie je webové. Postup práce je krátky:
\begin{enumerate}
	\item Otvorte aplikáciu v~prehliadači.
	\item Počkajte, kým stav v~hlavičke ukáže \emph{Online}.
	\item Vyberte model v~rozbaľovacom zozname.
	\item Zadajte otázku do vstupného poľa.
	\item Odošlite správu tlačidlom alebo klávesom Enter.
\end{enumerate}

Odpoveď sa zobrazí v~chatovom okne. Výstup má štruktúru ANSWER, KEY FACTS, RISK LEVEL a CONFIDENCE.

\section{Ovládacie prvky a tlačidlá}

V~nasledujúcich podsekciách sú popísané hlavné prvky rozhrania.

\subsection{Výber modelu}

Rozbaľovací zoznam v~hlavičke umožňuje zvoliť model. Zvolený model sa použije pre ďalšie odpovede.

\begin{figure}[ht]
	\centering
	\fbox{\parbox[c][4cm][c]{0.85\textwidth}{\centering
	\textit{\path{thesis.in.latex/figures/model_select.png}}}}
	\caption{Rozbaľovací zoznam pre výber modelu.}
	\label{fig:user-manual-model-select}
\end{figure}

\subsection{Indikátor stavu}

Indikátor zobrazuje dostupnosť služby. Stav \emph{Online} znamená, že API je pripravené.

\begin{figure}[ht]
	\centering
	\fbox{\parbox[c][4cm][c]{0.85\textwidth}{\centering
	\textit{\path{thesis.in.latex/figures/online.png}}}}
	\caption{Indikátor stavu služby.}
	\label{fig:user-manual-status}
\end{figure}

\subsection{Rýchle návrhy otázok}

Kliknutím na návrh sa text vloží do vstupného poľa. Ide o~rýchle ukážky tém.

\begin{figure}[ht]
	\centering
	\fbox{\parbox[c][4cm][c]{0.85\textwidth}{\centering
	\textit{\path{thesis.in.latex/figures/question_template.png}}}}
	\caption{Tlačidlá s~rýchlymi návrhmi otázok.}
	\label{fig:user-manual-suggestions}
\end{figure}

\subsection{Vstupné pole}

Do poľa sa píše otázka alebo inštrukcia. Enter odošle správu. Kombinácia Shift+Enter správu neodošle.

\begin{figure}[ht]
	\centering
	\fbox{\parbox[c][4cm][c]{0.85\textwidth}{\centering
	\textit{\path{thesis.in.latex/figures/question_field.png}}}}
	\caption{Vstupné pole pre zadanie otázky.}
	\label{fig:user-manual-input}
\end{figure}

\subsection{Tlačidlo Odoslať}

Tlačidlo so šípkou odošle aktuálnu otázku. Počas generovania je dočasne neaktívne.

\begin{figure}[ht]
	\centering
	\fbox{\parbox[c][4cm][c]{0.85\textwidth}{\centering
	\textit{\path{thesis.in.latex/figures/send_button.png}}}}
	\caption{Tlačidlo na odoslanie správy.}
	\label{fig:user-manual-send}
\end{figure}
